\section{Formalization}\label{sec:formalization}
\subsection{Representation}\label{sec:representation}
We write $\langle x_1,\ldots,x_n\rangle$ for a finite sequence and $\langle\rangle$
for the empty one; $|L|$ is the length of a sequence $L$, and $\#_v(L)$ the number
of occurrences of $v$ in it.

\begin{definition}[Schema]\label{def:schema}
A \emph{schema} $S$ consists of finite lists of declarations, each with an
identity and a simple name, that refer to each other by identity:
\begin{list}{}{\setlength{\leftmargin}{6mm}\setlength{\labelwidth}{5mm}%
 \setlength{\itemsep}{0pt}\setlength{\parsep}{0pt}\setlength{\topsep}{1pt}}
\item[(a)] packages, each with an optional parent package;
\item[(b)] classes, each with a package, direct superclasses and an abstract flag;
\item[(c)] properties $p$, each with an owner (a class or an association), a type
 (Boolean, integer, string, an enumeration or a class), bounds
 $\ell_p\in\mathbb N$ and $u_p\in\mathbb N\cup\{\ast\}$, where $\ast$ is
 unbounded, and flags $\mathit{unique}_p$, $\mathit{composite}_p$ and
 $\mathit{id}_p$ (identifier);
\item[(d)] associations, each with a package and two \emph{opposite ends} $(p,q)$;
\item[(e)] enumerations, each with a package, and literals, each with its enumeration.
\end{list}
\end{definition}

\begin{definition}[Snapshot]\label{def:snapshot}
A \emph{snapshot} $M=(O,\mathit{cls},R)$ consists of a finite list $O$ of object
identities, each with a class identity $\mathit{cls}(o)$, and a finite list $R$
of \emph{observation rows} $(o,p,L)$: an object identity, a property identity and
a sequence $L$ of values. A value is an object reference $\operatorname{ref}(o)$,
a Boolean, an integer, a string or an enumeration literal. We call $(S,M)$ a
\emph{Core} pair.
\end{definition}

\begin{definition}[Lookup]\label{def:lookup}
Let $(o,p,L_1),\ldots,(o,p,L_k)$ be the rows of $R$ for object $o$ and
property $p$, in their order in $R$. The \emph{lookup} of $p$ on $o$ is their
concatenation
\[
 L_M(o,p)=L_1L_2\cdots L_k,
\]
which is $\langle\rangle$ if $k=0$. We write $o.p$ for $L_M(o,p)$.
\end{definition}

\paragraph{Applicability.}
The \emph{ancestors} of a class are the class itself and the classes reached
through direct superclasses. A property $p$ \emph{applies} to class $c$, written
$p\in\mathit{app}(c)$, if an ancestor of $c$ declares $p$, or if $p$ is an
association-owned end whose opposite has an ancestor of $c$ as its type.
Following UML inheritance and EMOF property access
\cite{omg_uml_25,omg_mof_251}, inherited properties are collected by
declaration identity: a diamond yields one declaration, and same-named ones are
selected by qualified names, not overridden; e.g., \emph{RailwayElement.id}
applies to \emph{Switch}.

Figure~\ref{fig:running} shows our running snapshot $M_T$ as a UML object
diagram over the metamodel of Figure~\ref{fig:train}: route $r$ requires sensors
$s_A,s_B$ and follows switch position $\mathit{sp}$ of switch $\mathit{sw}$.
Each slot and each link end is one row. The two \emph{requires} links form the
single row $(r,\mathit{requires},\langle\operatorname{ref}(s_A),\operatorname{ref}(s_B)\rangle)$,
and the link between $r$ and $\mathit{sp}$ gives
$r.\mathit{follows}=\langle\operatorname{ref}(\mathit{sp})\rangle$ and
$\mathit{sp}.\mathit{route}=\langle\operatorname{ref}(r)\rangle$. Every other
applicable property, such as $r.\mathit{entry}$, has an empty row.

\begin{figure}[!t]
\centering
\def\objbox#1#2{\begin{tabular}{@{\hspace{3pt}}l@{\hspace{3pt}}}\underline{#1}\\[1.5pt]\hline #2\end{tabular}}
\resizebox{.8\linewidth}{!}{%
\begin{tikzpicture}[font=\footnotesize,
 obj/.style={draw,inner sep=0pt},
 nav/.style={-{Stealth[length=1.5mm]}},
 comp/.style={{Diamond[length=2.6mm,width=1.6mm]}-{Stealth[length=1.5mm]}},
 both/.style={{Stealth[length=1.5mm]}-{Stealth[length=1.5mm]}}]
\node[obj] (sA) at (0.4,0.65) {\objbox{$s_A$ : Sensor}{id = 404}};
\node[obj] (sB) at (0.4,-0.45) {\objbox{$s_B$ : Sensor}{id = 405}};
\node[obj] (r) at (4.0,0.1) {\objbox{$r$ : Route}{id = 101\\ active = true}};
\node[obj] (sp) at (9.3,0.85) {\objbox{$\mathit{sp}$ : SwitchPosition}{id = 303\\ position = STRAIGHT}};
\node[obj] (sw) at (9.3,-1.0) {\objbox{$\mathit{sw}$ : Switch}{id = 202\\ currentPosition = STRAIGHT}};
\draw[nav] (r.west) -- node[pos=.5,above,sloped]{requires} (sA.east);
\draw[nav] (r.west) -- node[pos=.5,below,sloped]{requires} (sB.east);
\draw[comp] (r.east) -- node[pos=.22,below,sloped]{route} node[pos=.72,above,sloped]{follows} (sp.west);
\draw[both] (sp.south) -- node[pos=.25,right]{positions} node[pos=.75,right]{target} (sw.north);
\end{tikzpicture}}
\caption{Running snapshot $M_T$ of the metamodel in Figure~\ref{fig:train}.}
\label{fig:running}
\end{figure}

With one row per pair, as in $M_T$, lookup returns that row's sequence, and an
empty row records absence: $r.\mathit{entry}=\langle\rangle$ differs from
$\langle\mathsf{false}\rangle$ or $\langle 0\rangle$. Lookup is also defined for
a missing row, giving $\langle\rangle$, and for repeated rows, which concatenate.

\subsection{Structural Conformance}\label{sec:conformance}
Definitions~\ref{def:wellformed} and~\ref{def:conformance} adapt structural
clauses of UML and MOF, as argued from their text and an audit of MOF~\S12.4's
32 EMOF restrictions. Below, we trace the conditions to the original
specifications and demonstrate them with counterexamples, i.e., variants of $M_T$
that each violate a single condition.

\begin{definition}[Well-formedness]\label{def:wellformed}
A schema $S$ is \emph{well formed}, $W(S)$, if:
\begin{list}{}{\setlength{\leftmargin}{9mm}\setlength{\labelwidth}{8mm}%
 \setlength{\itemsep}{1pt}\setlength{\parsep}{0pt}\setlength{\topsep}{2pt}}
\item[(W1)] \emph{Identities}: declarations of each kind have distinct identities.
\item[(W2)] \emph{Names}: every name is nonempty and consists of XML-valid characters.
\item[(W3)] \emph{Resolution}: every reference between declarations resolves to
 one of the expected kind, and an association owning a property has it as an end.
\item[(W4)] \emph{Acyclicity}: package nesting and generalization are acyclic.
\item[(W5)] \emph{Properties}: $\ell_p\leq u_p$; composite properties have class
 types; and a class has at most one identifier, declared or inherited.
\item[(W6)] \emph{Associations}: the two ends are distinct and class-typed, at
 most one is association-owned, and a class-owned end is owned by the type of its
 opposite; no property is an end of two associations; and the opposite of a composite
 end is not composite and has upper bound~1.
\end{list}
\end{definition}
W2, W5 and W6 include EMOF restrictions \cite{omg_mof_251}; W4 and the composite
conditions follow UML \cite{omg_uml_25}.

\begin{definition}[Conformance]\label{def:conformance}
A snapshot $M$ \emph{conforms} to a schema $S$, written $C(S,M)$, if $W(S)$
holds and, for all $o\in O$ and $p\in\mathit{app}(\mathit{cls}(o))$:
\begin{list}{}{\setlength{\leftmargin}{9mm}\setlength{\labelwidth}{8mm}%
 \setlength{\itemsep}{1pt}\setlength{\parsep}{0pt}\setlength{\topsep}{2pt}}
\item[(C1)] \emph{Objects}: identities are distinct, and $\mathit{cls}(o)$ is
 resolved and not abstract.
\item[(C2)] \emph{Rows}: every row names a stored object and a property that
 applies to its class, and the pair $(o,p)$ has exactly one row.
\item[(C3)] \emph{Typing}: every value in $L_M(o,p)$ fits the type of $p$:
 references to stored objects whose class has that type as an ancestor, literals
 of the enumeration, or XML-valid strings.
\item[(C4)] \emph{Multiplicity}: $\ell_p\leq|L_M(o,p)|\leq u_p$, and if
 $\mathit{unique}_p$, no value repeats in $L_M(o,p)$.
\item[(C5)] \emph{Opposites}: for each association $(p,q)$ and all $x,y\in O$,
 $\#_{\operatorname{ref}(y)}L_M(x,p)=\#_{\operatorname{ref}(x)}L_M(y,q)$.
\item[(C6)] \emph{Containment}: $o$ has at most one \emph{parent}, i.e., an
 $x\in O$ with $\operatorname{ref}(o)\in L_M(x,p')$ for a composite $p'$; at most
 one \emph{container end} $q$ of $o$, the opposite of a composite property, has
 $L_M(o,q)\neq\langle\rangle$; and no chain of parents leads from an object back
 to itself.
\end{list}
\end{definition}

\paragraph{Objects, rows and typing (C1--C3).}
C1 follows UML: abstract classes have no direct instances \cite{omg_uml_25}.
C2 demands the rows of Section~\ref{sec:representation}, recording MOF's null
or empty list \cite{omg_mof_251} as $\langle\rangle$, so it rejects missing and
repeated rows; deleting $\mathit{sw}$'s \emph{id} row violates it alone. C3
follows UML type conformance and MOF's XML Schema datatypes.

\paragraph{Multiplicity (C4).}
Following UML \cite{omg_uml_25}, bounds count occurrences, and uniqueness
separately excludes repeats. Dropping $s_B$ from $M_T$ leaves
$|r.\mathit{requires}|=1<2$, which violates C4 alone. A $0..0$ bound means that
a property stays empty; MOF requires positive upper bounds only before
reflective \texttt{create} \cite[\S9.3.3]{omg_mof_251}, which our subset omits.

\paragraph{Opposites (C5).}
Without link identities, an association is observed at both ends, and C5 makes
the ends agree occurrence by occurrence. Clearing $\mathit{sp}.\mathit{route}$
in $M_T$ gives $0\neq1$ against $r.\mathit{follows}$: a link recorded at one end
only. For unique ends, C5 is reciprocal membership; with mixed uniqueness, it is
stricter than UML \cite[\S11.5.3.1]{omg_uml_25}. Were the unique
\emph{Route.follows} nonunique, UML would permit two links between $r$ and
$\mathit{sp}$, seen as
$\langle\operatorname{ref}(\mathit{sp}),\operatorname{ref}(\mathit{sp})\rangle$
at \emph{follows} but $\langle\operatorname{ref}(r)\rangle$ at the unique
\emph{route} end; C5 rejects this ($2\neq1$), a restriction of our subset.

\paragraph{Containment (C6).}
C6 follows MOF \cite[\S12.5.5]{omg_mof_251}. Parents form a set, so repeated
occurrences add none: a region listing
$\langle\operatorname{ref}(\mathit{seg}),\operatorname{ref}(\mathit{seg})\rangle$
gives $\mathit{seg}$ one parent. Only the uniqueness of \emph{Region.elements}
(C4) rejects this list; with uniqueness disabled, it is accepted. A segment in the
\emph{elements} of two regions violates C6 alone.

\subsection{Verified Checking and Translation}\label{sec:guarantees}
We prove that the checker accepts exactly the conforming pairs and that
elaboration preserves meaning; Section~\ref{sec:tool} describes the Lean
mechanization.

\begin{theorem}[Checker correctness]\label{thm:checker}
For every schema $S$ and snapshot $M$, the checker's Boolean verdict satisfies
\begin{equation}
 \operatorname{check}(S,M)=\mathsf{true}
 \quad\Longleftrightarrow\quad C(S,M).
\label{eq:checker}
\end{equation}
\end{theorem}
\begin{proofsketch}
Apart from reachability, each condition of Definitions~\ref{def:wellformed}
and~\ref{def:conformance} quantifies over the finite lists of $S$ and $M$, e.g.,
C4 over $O$ and $\mathit{app}(\mathit{cls}(o))$, and the checker evaluates it
element by element, which is exact. For the ancestors of a class $c$, it
computes, as breadth-first search does, $X_0=\{c\}$ and
$X_{k+1}=X_k\cup\mathit{sup}(X_k)$, where $\mathit{sup}(X)$ holds the direct
superclasses of the classes in $X$; for \emph{Switch}, $X_1$ adds
\emph{TrackElement}, $X_2$ adds \emph{RailwayElement}, and $X_3=X_2$. Under W3,
$X_0\subseteq X_1\subseteq\cdots$ lie among the $n$ classes of $S$, so $X_{n+1}=X_n$:
$X_n$ is closed under $\mathit{sup}$, hence contains every ancestor, and it
contains only ancestors. Package nesting (W4) and parent chains (C6) are alike;
there, W3, C2 and C3 keep links within the lists. A pair violating these does not conform and is rejected, as the
checker evaluates them too.
\end{proofsketch}

\paragraph{From source documents to Core.}
Theorem~\ref{thm:checker} concerns Core pairs, but users write text, which a
parser turns into a source document that uses names instead of identities.
Translating it into Core must preserve its meaning: dropping a repeated value
would hide a uniqueness violation (C4), and merging two names would create one.
We therefore define the translation and, independently of it, what a source
document means.

\begin{definition}[Source document]\label{def:document}
A \emph{source document} $d$ has the components of Definitions~\ref{def:schema}
and~\ref{def:snapshot}, with names in place of identities: each declaration has
a \emph{qualified name}, a sequence of strings such as
$\mathit{Route}{::}\mathit{requires}$, besides its simple name; each object has an
\emph{object name}, such as $r$; and references and rows use these names. $L_d$ is
the lookup of $d$. We call $d$ \emph{well named}, i.e., its names can serve as identities, if:
\begin{list}{}{\setlength{\leftmargin}{9mm}\setlength{\labelwidth}{8mm}%
 \setlength{\itemsep}{1pt}\setlength{\parsep}{0pt}\setlength{\topsep}{2pt}}
\item[(N1)] every qualified and object name is nonempty and has nonempty components;
\item[(N2)] no two declarations, and no two objects, have the same name;
\item[(N3)] each declaration's qualified name extends its owner's by one
 component, or has one component if the declaration has no owner;
\item[(N4)] every name used in $d$ denotes a declaration or object of the expected kind.
\end{list}
We call $d$ \emph{lexically admissible} if the components of its qualified and
object names consist of XML-valid characters.
\end{definition}

\begin{definition}[Elaboration]\label{def:elaboration}
\emph{Elaboration} $E$ translates a source document $d$ into a Core pair. It fails
unless $d$ is well named; otherwise $E(d)=\mathsf{ok}(\rho(d))$. Here $\rho$
numbers the declarations of each kind and the objects by their positions in $d$,
and $\rho(d)$ replaces every qualified or object name $n$ by $\rho(n)$, keeping
simple names and the order and repetition of values.
\end{definition}
Elaborating the source document of $M_T$, whose objects are named $r$, $s_A$,
$s_B$, $\mathit{sp}$ and $\mathit{sw}$, yields $M_T$ up to the choice of
identities; naming two objects $r$ makes elaboration fail by (N2).

\begin{definition}[Source meaning]\label{def:meaning}
The \emph{meaning} of a source document $d$ is the condition $\mathcal S(d)$ that
$d$ is well named and lexically admissible and that $W$ and C1--C6 hold on $d$
with names in place of identities.
\end{definition}
$\mathcal S(d)$ is stated on $d$ alone, without $E$ or $\rho$. Elaboration checks every name condition of $\mathcal S(d)$ except lexical
admissibility, which a Core pair cannot confirm, since it keeps no qualified or
object names.

\begin{theorem}[Elaboration adequacy]\label{thm:adequacy}
Let $d$ be a lexically admissible source document. (i)~If $\mathcal S(d)$, then
$E(d)$ succeeds. (ii)~If $E(d)=\mathsf{ok}(S,M)$, then $\mathcal S(d)\Leftrightarrow C(S,M)$.
\end{theorem}
Part~(i) says that elaboration rejects only documents whose meaning fails, and
part~(ii) that a successful elaboration neither hides nor introduces violations.
\begin{proofsketch}
(i)~Elaboration fails only if $d$ is not well named, which $\mathcal S(d)$ rules
out. (ii)~By (N2), $\rho$ is a bijection from the names of $d$ onto the identities of
$E(d)=\rho(d)$, which keeps every row and the order and repetition of its
values; hence $L_{E(d)}(\rho o,\rho p)=\rho(L_d(o,p))$, and $\rho$ is an
isomorphism from $d$ onto $E(d)$. $W$ and C1--C6 refer to identities only
through equality, lookup and references, which isomorphisms preserve, so each
holds on $d$ exactly when it holds on $E(d)$. The faulty translations above are
not isomorphisms: dropping a repeat changes a lookup, and merging two names
breaks injectivity. For $C(S,M)\Rightarrow\mathcal S(d)$, only the name conditions
remain: $d$ is well named because $E(d)$ succeeded, and lexically admissible by
hypothesis.
\end{proofsketch}
With Theorem~\ref{thm:checker}, every lexically admissible source document $d$
with $E(d)=\mathsf{ok}(S,M)$ thus satisfies
\begin{equation}
 \mathcal S(d)
 \ \overset{\text{Thm.~\ref{thm:adequacy}}}{\Longleftrightarrow}\
 C(S,M)
 \ \overset{\text{Thm.~\ref{thm:checker}}}{\Longleftrightarrow}\
 \operatorname{check}(S,M)=\mathsf{true},
\label{eq:adequacy}
\end{equation}
and by~(i), $\mathcal S(d)$ fails whenever $E(d)$ fails. Checking the elaborated
pair thus answers exactly whether a source document satisfies its meaning.
